\chapter{1982年全国卷（理）}

\section{填空题}

\begin{problem}

填表：

\begin{center}
\renewcommand{\arraystretch}{1.2}
\setlength{\tabcolsep}{6pt}
\begin{tabular}{|c|c|c|}
\hline
 & 函数 & 使函数有意义的 \(x\) 的实数范围 \\
\hline
1 & \(y=\sqrt{-x^2}\) &\fillinblank{} \\
\hline
2 & \(y=\sqrt{(-x)^2}\) &\fillinblank{} \\
\hline
3 & \(y=\arcsin(\sin x)\) &\fillinblank{} \\
\hline
4 & \(y=\sin(\arcsin x)\) &\fillinblank{} \\
\hline
5 & \(y=10^{\lg x}\) &\fillinblank{} \\
\hline
6 & \(y=\lg 10^x\) &\fillinblank{} \\
\hline
\end{tabular}
\end{center}
\end{problem}

\begin{answer}
依次为 \(\{0\}\)，\(\R\)，\(\R\)，\([-1,1]\)，\((0,+\infty)\)，\(\R\).
\end{answer}

\begin{solution}
\begin{enumerate}
\item 根式有意义要求被开方数非负，即 \(-x^2\geqslant0\)，所以 \(x^2\leqslant0\)，从而 \(x=0\)，定义域为 \(\{0\}\).
\item 因为 \((-x)^2=x^2\geqslant0\) 对任意实数 \(x\) 都成立，所以定义域为 \(\R\).
\item 对任意 \(x\in\R\)，都有 \(-1\leqslant\sin x\leqslant1\)，而 \(\arcsin t\) 在 \([-1,1]\) 上有意义，所以定义域为 \(\R\).
\item \(\arcsin x\) 有意义的充要条件是 \(-1\leqslant x\leqslant1\)，故定义域为 \([-1,1]\).
\item \(\lg x\) 有意义的充要条件是 \(x>0\)，故定义域为 \((0,+\infty)\).
\item 对任意实数 \(x\)，都有 \(10^x>0\)，所以 \(\lg 10^x\) 对任意实数 \(x\) 有意义，定义域为 \(\R\).
\end{enumerate}
\end{solution}

\section{解答题}

\begin{problem}

\begin{enumerate}
\item 求 \((-1+\i)^{20}\) 展开式中第 \(15\) 项的数值；
\item 求 \(y=\cos^2 \frac{x}{3}\) 的导数.
\end{enumerate}
\end{problem}

\begin{solution}
\begin{enumerate}
\item 由二项式定理，展开式的第 \(15\) 项对应 \(i\) 的幂为 \(14\)，因此
\[
T_{15}=\mathrm C_{20}^{14}(-1)^{20-14}\i^{14}=\mathrm C_{20}^{14}\i^{14}=-\mathrm C_{20}^{14}=-38760
\]
这里用到了 \(\i^{14}=\i^{12}\i^2=-1\).
\item 令 \(u=\cos\dfrac{x}{3}\)，则 \(y=u^2\)，且
\[
u'=-\dfrac{1}{3}\sin\dfrac{x}{3}
\]
所以
\[
y'=2u u'=-\dfrac{2}{3}\cos\dfrac{x}{3}\sin\dfrac{x}{3}=-\dfrac{1}{3}\sin\dfrac{2x}{3}
\]
\end{enumerate}
\end{solution}

\begin{problem}

在平面直角坐标系内，下列方程表示什么曲线？画出它们的图形.

\begin{enumerate}
\item \(\begin{vmatrix} 2x & 1 & 1 \\ -3y & 2 & 3 \\ 6 & 3 & 4 \end{vmatrix}=0\)；
\item \(\begin{cases}
x=1+\cos \varphi,\\
y=2\sin \varphi
\end{cases}\).
\end{enumerate}
\end{problem}

\begin{solution}
\begin{enumerate}
\item 按第一行展开行列式，得
\[
\begin{vmatrix} 2x & 1 & 1 \\ -3y & 2 & 3 \\ 6 & 3 & 4 \end{vmatrix}
=2x(8-9)-\bigl((-3y)\cdot4-3\cdot6\bigr)+\bigl((-3y)\cdot3-2\cdot6\bigr)
=-2x+3y+6
\]
因此原方程等价于 \(2x-3y-6=0\)，即 \(y=\dfrac{2}{3}x-2\)，表示一条直线. 它与 \(x\) 轴、\(y\) 轴分别交于 \((3,0)\)、\((0,-2)\)，连接这两点即可作出图形.
\item 由参数方程有 \(x-1=\cos\varphi\)、\(y=2\sin\varphi\)，于是
\[
\left(x-1\right)^2+\dfrac{y^2}{4}=\cos^2\varphi+\sin^2\varphi=1
\]
当 \(\varphi\) 取遍实数时，这条曲线的全部点都被取得，故它表示以 \((1,0)\) 为中心、半轴长分别为 \(1\) 和 \(2\) 的椭圆. 其顶点为 \((0,0)\)、\((2,0)\)、\((1,2)\)、\((1,-2)\)，据此即可作图.
\end{enumerate}
\end{solution}

\begin{problem}

已知圆锥体的底面半径为 \(R\)，高为 \(H\)，求内接于这个圆锥体并且体积最大的圆柱体的高 \(h\)（如图）.

\bitmapfigure[float=inline,max width=0.44\linewidth]{img/1982/national_science/q4_fig1.png}
\FloatBarrier
\end{problem}

\begin{solution}
设圆柱的底面半径为 \(r\)，高为 \(h\)，其中 \(0<h<H\). 由圆锥的轴截面中相似三角形关系，圆柱顶面到圆锥顶点的距离为 \(H-h\)，故
\[
\dfrac{r}{R}=\dfrac{H-h}{H}\Longrightarrow r=\dfrac{R(H-h)}{H}
\]
圆柱体积为
\[
V(h)=\pi r^2h=\dfrac{\pi R^2}{H^2}h(H-h)^2
\]
在 \(0<h<H\) 内求导，得
\[
V'(h)=\dfrac{\pi R^2}{H^2}(H-h)(H-3h)
\]
因此 \(0<h<\dfrac{H}{3}\) 时 \(V'(h)>0\)，而 \(\dfrac{H}{3}<h<H\) 时 \(V'(h)<0\). 体积在 \(h=\dfrac{H}{3}\) 处由增转减，且两端的极限体积均为 \(0\)，所以最大体积对应的圆柱高为
\[
h=\dfrac{H}{3}
\]
\end{solution}

\begin{problem}

设 \(0<x<1\)，\(a>0\)，\(a\neq 1\)，比较 \(\left|\log_a(1-x)\right|\) 与 \(\left|\log_a(1+x)\right|\) 的大小. （要写出比较过程）
\end{problem}

\begin{solution}
因为 \(a>0\) 且 \(a\neq1\)，由换底公式
\[
\left|\log_a t\right|=\dfrac{|\ln t|}{|\ln a|}
\]
又因为 \(0<x<1\)，所以 \(0<1-x<1<1+x\)，从而
\[
\left|\log_a(1-x)\right|=\dfrac{-\ln(1-x)}{|\ln a|}
\]
以及
\[
\left|\log_a(1+x)\right|=\dfrac{\ln(1+x)}{|\ln a|}
\]
比较分子，注意到 \(0<1-x^2<1\)，于是 \(\ln(1-x^2)<0\)，即
\[
-\ln(1-x)-\ln(1+x)=-\ln(1-x^2)>0
\]
因此
\[
\left|\log_a(1-x)\right|>\left|\log_a(1+x)\right|
\]
\end{solution}

\begin{problem}

如图：已知锐角 \(\angle AOB=2\alpha\) 内有动点 \(P\)，\(PM\perp OA\)，\(PN\perp OB\)，且四边形 \(PMON\) 的面积等于常数 \(c^2\). 今以 \(O\) 为极点，\(\angle AOB\) 的角平分线 \(OX\) 为极轴，求动点 \(P\) 的轨迹的极坐标方程，并说明它表示什么曲线.

\bitmapfigure[float=inline,max width=0.46\linewidth]{img/1982/national_science/q6_fig1.png}
\end{problem}

\begin{solution}
设 \(OP=\rho\)，且以角平分线 \(OX\) 为极轴，记有向角 \(\angle XOP=\theta\). 由于 \(P\) 在角 \(AOB\) 内，\(-\alpha<\theta<\alpha\). 在直角三角形 \(OMP\)、\(ONP\) 中，分别有 \(OM=\rho\cos(\alpha-\theta)\)，\(PM=\rho\sin(\alpha-\theta)\).

同理，\(ON=\rho\cos(\alpha+\theta)\)，\(PN=\rho\sin(\alpha+\theta)\).
对角线 \(OP\) 将四边形 \(PMON\) 分成这两个直角三角形，所以
\[
c^2=\dfrac12OM\cdot PM+\dfrac12ON\cdot PN
=\dfrac{\rho^2}{4}\left[\sin(2\alpha-2\theta)+\sin(2\alpha+2\theta)\right]
\]
再由和差化积公式，得
\[
c^2=\dfrac{\rho^2}{2}\sin2\alpha\cos2\theta
\]
故动点轨迹的极坐标方程为
\[
\rho^2\cos2\theta=\dfrac{2c^2}{\sin2\alpha}
\]
由 \(x=\rho\cos\theta\)、\(y=\rho\sin\theta\) 以及 \(\rho^2\cos2\theta=x^2-y^2\)，可化为
\[
x^2-y^2=\dfrac{2c^2}{\sin2\alpha}
\]
右端为正，且 \(P\) 在角 \(AOB\) 内时 \(x>0\)，所以它表示双曲线 \(x^2-y^2=\dfrac{2c^2}{\sin2\alpha}\) 的右支在角 \(AOB\) 内的一段.
\end{solution}

\begin{problem}

已知空间四边形 \(ABCD\) 中 \(AB=BC\)，\(CD=DA\)，\(M\)，\(N\)，\(P\)，\(Q\) 分别是边 \(AB\)，\(BC\)，\(CD\)，\(DA\) 的中点（如图）. 求证：\(MNPQ\) 是一个矩形.

\begin{center}
\bitmapfigure[float=inline,max width=0.36\linewidth]{img/1982/national_science/q7_fig1.png}
\end{center}
\FloatBarrier
\end{problem}

\begin{solution}
设点 \(A\)，\(B\)，\(C\)，\(D\) 的位置向量分别为 \(\bs{a}\)，\(\bs{b}\)，\(\bs{c}\)，\(\bs{d}\). 由中点公式，
\[
\overrightarrow{MN}=\dfrac12(\bs{c}-\bs{a})=\dfrac12\overrightarrow{AC}
\]
又有
\[
\overrightarrow{QP}=\dfrac12(\bs{c}-\bs{a})
\]
以及
\[
\overrightarrow{NP}=\dfrac12(\bs{d}-\bs{b})=\overrightarrow{MQ}
\]
所以 \(MNPQ\) 是平行四边形，且 \(MN\parallel AC\)、\(NP\parallel BD\).

令 \(\bs{u}=\bs{c}-\bs{a}\)，\(\bs{v}=\bs{b}-\bs{a}\)，\(\bs{w}=\bs{d}-\bs{a}\). 由 \(AB=BC\)，有
\[
|\bs{v}|^2=|\bs{u}-\bs{v}|^2
\]
从而 \(\bs{u}\cdot\bs{v}=\dfrac12|\bs{u}|^2\). 由 \(CD=DA\)，有
\[
|\bs{u}-\bs{w}|^2=|\bs{w}|^2
\]
从而 \(\bs{u}\cdot\bs{w}=\dfrac12|\bs{u}|^2\). 因此
\[
\overrightarrow{AC}\cdot\overrightarrow{BD}=\bs{u}\cdot(\bs{w}-\bs{v})=0
\]
即 \(AC\perp BD\). 于是 \(MN\perp NP\)，平行四边形 \(MNPQ\) 有一个直角，故 \(MNPQ\) 是矩形.
\end{solution}

\begin{problem}

抛物线 \(y^2=2px\) 的内接三角形有两边与抛物线 \(x^2=2qy\) 相切，证明这个三角形的第三边也与 \(x^2=2qy\) 相切.
\end{problem}

\begin{solution}
由于题中两式均表示非退化抛物线，\(p\)、\(q\) 均不为零. 将 \(y^2=2px\) 上参数为 \(t\) 的点写成
\[
T(t)=\left(\dfrac{p}{2}t^2,pt\right)
\]
参数为 \(u\)，\(v\) 的两点所确定的弦方程为
\[
2x-(u+v)y+puv=0
\]
若这条弦与 \(x^2=2qy\) 相切，先证明 \(u+v\neq0\). 若 \(u+v=0\)，弦方程为 \(2x-pu^2=0\)，它是竖直直线；而 \(q\neq0\) 时，抛物线 \(x^2=2qy\) 在任一点 \(\left(x_0,y_0\right)\) 处由隐式求导得 \(\dfrac{dy}{dx}=\dfrac{x_0}{q}\)，故切线方程为 \(2x_0x-2qy-x_0^2=0\)，其斜率有限，不可能是竖直直线，矛盾. 因此可写为
\[
y=\dfrac{2}{u+v}x+\dfrac{puv}{u+v}
\]
将其与 \(x^2=2qy\) 联立，所得关于 \(x\) 的二次方程必须有重根，故
\[
\left(\dfrac{4q}{u+v}\right)^2+\dfrac{8qpuv}{u+v}=0
\]
即
\[
puv(u+v)+2q=0
\]

设三角形三个顶点的参数为 \(u\)，\(v\)，\(w\)，不妨设边 \(T(u)T(v)\)、\(T(v)T(w)\) 与第二条抛物线相切. 由这个相切条件，分别有
\[
puv(u+v)+2q=0
\]
和
\[
pvw(v+w)+2q=0
\]
由于 \(q\neq0\)，有 \(v\neq0\)，两式相减得
\[
u(u+v)=w(v+w)
\]
即
\[
(u-w)(u+v+w)=0
\]
三角形非退化，所以 \(u\neq w\)，从而 \(u+v+w=0\). 又因 \(v\neq0\)，有 \(u+w=-v\neq0\)，因此
\[
puw(u+w)+2q=-puvw+2q=puv(u+v)+2q=0
\]
再由同一相切判定，边 \(T(u)T(w)\) 也与 \(x^2=2qy\) 相切，证毕.
\end{solution}

\section{附加题}

\begin{problem}

已知数列 \(a_1\)，\(a_2\)，\(\cdots\)，\(a_n\)，\(\cdots\) 和数列 \(b_1\)，\(b_2\)，\(\cdots\)，\(b_n\)，\(\cdots\)，其中 \(a_1=p\)，\(b_1=q\)，\(a_n=pa_{n-1}\)，\(b_n=qa_{n-1}+rb_{n-1}\)（\(n\geqslant 2\)），\(p\)，\(q\)，\(r\) 是已知常数，且 \(q\neq 0\)，\(p>r>0\).

\begin{enumerate}
\item 用 \(p\)，\(q\)，\(r\)，\(n\) 表示 \(b_n\)，并用数学归纳法加以证明；
\item 求 \(\displaystyle\lim_{n\to\infty}\frac{b_n}{\sqrt{a_n^2+b_n^2}}\).
\end{enumerate}
\end{problem}

\begin{solution}
\begin{enumerate}
\item 先由 \(a_1=p\) 及递推式归纳得到
\[
a_n=p^n
\]
确实，\(a_1=p=p^1\)；若 \(a_{n-1}=p^{n-1}\)，则 \(a_n=pa_{n-1}=p^n\).

据此猜想
\[
b_n=\dfrac{q\left(p^n-r^n\right)}{p-r}
\]
当 \(n=1\) 时，右端为 \(q(p-r)/(p-r)=q=b_1\)，命题成立. 假设对某个 \(n-1\geqslant1\) 成立，则
\[
b_n=qa_{n-1}+rb_{n-1}=qp^{n-1}+r\cdot\dfrac{q\left(p^{n-1}-r^{n-1}\right)}{p-r}
\]
\[
b_n=\dfrac{q\left((p-r)p^{n-1}+rp^{n-1}-r^n\right)}{p-r}
=\dfrac{q\left(p^n-r^n\right)}{p-r}
\]
所以由数学归纳法，对所有正整数 \(n\) 都有
\[
b_n=\dfrac{q\left(p^n-r^n\right)}{p-r}
\]

\item 因为 \(p>r>0\)，所以 \(p^n>0\)，并且
\[
\dfrac{b_n}{p^n}=\dfrac{q}{p-r}\left(1-\left(\dfrac{r}{p}\right)^n\right)\longrightarrow\dfrac{q}{p-r}
\]
于是
\[
\frac{b_n}{\sqrt{a_n^2+b_n^2}}
=\frac{b_n/p^n}{\sqrt{1+(b_n/p^n)^2}}
\]
因此
\[
\lim_{n\to\infty}\frac{b_n}{\sqrt{a_n^2+b_n^2}}
=\frac{q/(p-r)}{\sqrt{1+q^2/(p-r)^2}}
=\dfrac{q}{\sqrt{(p-r)^2+q^2}}
\]
\end{enumerate}
\end{solution}
